\chapter*{Послесловие: будущее обучения с подкреплением}
\addcontentsline{toc}{chapter}{Послесловие}

Если вы дочитали до этого места, значит, вы проследовали путь от игрока у игрового автомата до алгоритмов, определяющих то, как передовые системы ИИ общаются с людьми. Этот путь оказывается удивительно прямым. Принцип оптимальности, сформулированный Ричардом Беллманом в 1950-х годах; уравнения временных разностей, разработанные Саттоном и Барто в 1980-х; теорема о градиенте стратегии, доказанная на рубеже тысячелетий\,---\,всё это не музейные экспонаты. Это живой механизм, работающий внутри систем, с которыми сотни миллионов людей взаимодействуют ежедневно. Редко когда разрыв между математической абстракцией и практической инженерией закрывался столь стремительно, и стоит остановиться, в конце этого пути, чтобы поразмышлять о том, куда движется область.

\medskip

\textbf{Обучение с подкреплением в центре ИИ.}\quad
На протяжении большей части своей истории обучение с подкреплением занимало важный, но несколько периферийный угол в ландшафте машинного обучения. На практике доминировало обучение с учителем; RL было предметом красивой теории и впечатляющих демонстраций\,---\,Atari, Go, локомоция роботов\,---\,однако его влияние на развёрнутые системы оставалось ограниченным. Это кардинально изменилось.

Сегодня RL занимает центральное место по меньшей мере в четырёх наиболее активных направлениях исследований в области ИИ. В сфере \emph{языка и рассуждения} выравнивание больших языковых моделей на основе обратной связи от людей и обучение моделей, способных последовательно рассуждать в рамках развёрнутых цепочек мыслей, оба зависят от оптимизации стратегий. В области \emph{агентных систем} модели, которые планируют на несколько шагов вперёд, используют инструменты, просматривают веб и пишут и выполняют код, являются по существу RL-агентами, действующими в структурированных средах с отложенными вознаграждениями. В \emph{мультимодальном ИИ} распространение обучения в стиле RLHF на зрение, аудио и видео~--- активно и стремительно развивающаяся область. В \emph{робототехнике и воплощённом интеллекте} конвейер sim-to-real, переносящий стратегии, обученные в симуляции, на физическое оборудование, остаётся одним из наиболее перспективных путей к созданию способных физических агентов. Во всех этих областях концептуальный словарь~--- именно тот, которым вы теперь владеете: состояния, действия, вознаграждения, стратегии, функции ценности и алгоритмы, их оптимизирующие.

\medskip

\textbf{Сближение двух традиций.}\quad
Одним из наиболее поразительных интеллектуальных достижений прошедшего десятилетия стало открытие, что классическое RL и обучение больших языковых моделей не просто аналогичны, но глубоко едины. Задача дообучения языковой модели является марковским процессом принятия решений: состояние~--- это история разговора, действие~--- следующий токен, стратегия~--- языковая модель, а вознаграждение~--- сигнал\,---\,предоставленный людьми или полученный от модели\,---\,о качестве ответа.

Это осознание означает, что каждый результат данной книги применим в принципе к языковым моделям. Анализ смещения-дисперсии методов Монте-Карло в сравнении с методами временных разностей вновь возникает в дискуссиях о приписывании заслуг в длинных цепочках рассуждений. Теорема о градиенте стратегии лежит в основе оптимизации вознаграждения в стиле REINFORCE. Архитектура «актор-критик» является прямым предком конвейеров RLHF на основе PPO, применяемых на практике. Уравнения Беллмана, записанные в эпоху исследования операций, теперь управляют обучением систем, которые сочиняют поэзию и доказывают теоремы.

Это сближение~--- не случайность. Оно отражает универсальность основы последовательного принятия решений и предполагает, что дальнейший прогресс в выравнивании языковых моделей будет частично связан с более глубоким погружением в теоретический инструментарий, накопленный RL за семьдесят лет.

\medskip

\textbf{Открытые вызовы.}\quad
Было бы нечестно завершать на ноте чистого триумфа. Обучение с подкреплением, при всех недавних успехах, сталкивается с серьёзными нерешёнными проблемами, и интеллектуальная честность требует их назвать.

\emph{Выборочная эффективность} остаётся фундаментальным ограничением. Люди научаются играть в большинство игр достаточно хорошо за несколько часов практики; лучшим RL-агентам требуются миллионы взаимодействий со средой. Преодоление этого разрыва~--- посредством более эффективного исследования, структурированных моделей мира или трансферного обучения~--- является одним из наиболее глубоких открытых вопросов в области.

\emph{Безопасность и робастность} становятся всё более насущными. Взлом вознаграждения, когда агент находит непредусмотренные способы максимизировать некорректно заданный сигнал вознаграждения,~--- это не просто теоретический курьёз; он наблюдался в развёрнутых системах. Обеспечение надёжного поведения RL-агентов при изменении среды, несовершенстве модели вознаграждения или наличии противников~--- проблема как теоретической, так и практической важности.

\emph{Координация множества агентов} вводит новый уровень сложности, с которым теория одного агента не справляется. Системы, состоящие из множества взаимодействующих агентов\,---\,будь то кооперирующие роботы или конкурирующие участники рынка\,---\,демонстрируют эмерджентную динамику, которая остаётся плохо изученной и трудно управляемой.

\emph{Масштабируемый надзор}~--- пожалуй, наиболее самобытный современный вызов. По мере того как системы ИИ становятся более способными, люди-оценщики всё меньше способны надёжно судить о качестве их результатов. Как поддерживать значимый контроль со стороны человека над системой, рассуждения которой уже могут превосходить возможности оценщика по верификации? Конституциональный ИИ, дебаты и рекурсивное моделирование вознаграждения~--- ранние попытки ответить на этот вопрос, однако удовлетворительное решение остаётся недостижимым.

\emph{Спецификация вознаграждения} лежит в основе всего вышесказанного. Проблема выражения человеческих ценностей, целей и предпочтений в виде функции вознаграждения, допускающей оптимизацию без нежелательных побочных эффектов, является по своей сути философской задачей в не меньшей мере, чем технической. Прогресс здесь потребует сотрудничества между сообществом RL, исследователями в области выравнивания, этиками и экспертами предметных областей.

\medskip

\textbf{Личное замечание.}\quad
Написание учебника~--- это акт оптимизма. Он предполагает, что идеи, содержащиеся в нём, достойны передачи, что время читателя, потраченное на их изучение, потрачено хорошо, и что область станет лучше от того, что в ней появится больше людей, понимающих её глубоко. Мы искренне придерживаемся всех трёх этих убеждений.

Благодарим вас за время и внимание, вложенные в этот материал. Изучение обучения с подкреплением нетривиально: математика требовательна, интуиция временами контринтуитивна, а разрыв между теорией и работающей реализацией может смирять гордость. Если вы прошли этот путь, вы лучше большинства практиков подготовлены к тому, чтобы рассуждать вдумчиво о системах, которые будут определять облик наступающего десятилетия.

Мы надеемся, что вы используете эти знания не просто для создания более способных систем, но для создания лучших\,---\,систем, которые безопасны, интерпретируемы и подлинно полезны. Область нуждается в людях, которые серьёзно относятся и к возможностям, и к ответственности, и мы верим, что теперь вы готовы внести вклад в оба направления.

\medskip

\textbf{Заключительная мысль.}\quad
В предисловии мы описали обучение с подкреплением как изучение того, как агент учится действовать правильно в мире, который он не вполне понимает, руководствуясь разреженным и порой вводящим в заблуждение сигналом вознаграждения. Это описание применимо с небольшими поправками к самому человеческому существованию. Возможно, именно поэтому эта концептуальная схема находит столь глубокий отклик и почему лучшие исследователи в области движимы чем-то большим, чем технический интерес.

Обучение с подкреплением в самых высоких своих амбициях~--- это теория о том, как разум может совершенствоваться через опыт. Если нам удастся всё сделать правильно~--- математику, инженерию и ценности, заложенные в сигналы вознаграждения, которые мы выбираем,~--- оно может оказаться одним из наиболее важных интеллектуальных вкладов нашей эпохи. Работа далека от завершения. Мы призываем вас продолжать там, где эта книга заканчивается.

\bigskip

\hfill\textit{Москва, 2026}
